Ce projet de fin d'études présente la conception et la réalisation de \textbf{PharmacieConnect}, une
plateforme web et mobile dédiée à la recherche de pharmacies, à la consultation des pharmacies de garde
et à l'administration des données pharmaceutiques. L'idée du projet part d'un besoin concret : les
informations sur les pharmacies de garde sont souvent éparpillées, difficiles à mettre à jour et pas
toujours pratiques à consulter depuis un téléphone.

Pour répondre à ce besoin, la solution a été organisée autour de trois parties. Le backend, développé avec
\textbf{FastAPI}, centralise les données, expose des API REST et applique les règles métier. Il assure aussi
la sécurité des accès grâce à une authentification par jetons JWT. Le tableau de bord administrateur,
réalisé avec \textbf{React} et \textbf{Vite}, permet de gérer les pharmacies, les plannings de garde, les
médicaments, les utilisateurs autorisés, les imports CSV, les indicateurs et les journaux d'audit.
L'application mobile, développée avec \textbf{Expo} et \textbf{React Native}, permet au citoyen de rechercher
une pharmacie, de consulter les gardes par date, d'utiliser la géolocalisation et d'accéder au catalogue
des médicaments.

Le travail réalisé couvre également plusieurs aspects de génie logiciel : modélisation des données avec
SQLAlchemy, validation des entrées, contrôle des rôles, journalisation des actions sensibles, tests
backend et séparation claire entre les interfaces clientes, les services métier et la couche de persistance.

\textbf{Mots-clés :} pharmacie de garde, FastAPI, React, Vite, Expo, React Native, JWT, API REST,
SQLAlchemy, tableau de bord, application mobile.
